\documentclass[11pt]{article}
\usepackage[T1]{fontenc}
\usepackage[utf8]{inputenc}
\usepackage{lmodern}
\usepackage{microtype}
\usepackage[a4paper,margin=1in]{geometry}
\usepackage{amsmath,amssymb,amsthm,mathtools}
\usepackage{booktabs,longtable,array}
\usepackage{enumitem}
\usepackage{xcolor}
\usepackage{hyperref}
\usepackage[nameinlink,noabbrev]{cleveref}
\usepackage{fancyhdr}
\usepackage{url}
\usepackage{graphicx}

\hypersetup{
  colorlinks=true,
  linkcolor=blue!45!black,
  citecolor=green!35!black,
  urlcolor=blue!55!black,
  pdftitle={Constructive Central Limit Theorems in iPRA},
  pdfauthor={OpenAI, prepared for Vladimir Reshetnikov}
}
\setlength{\emergencystretch}{3em}
\setlength{\headheight}{14pt}
\setlist{nosep,leftmargin=2em}
\pagestyle{fancy}
\fancyhf{}
\fancyhead[L]{Constructive central limit theorems in iPRA}
\fancyhead[R]{\thepage}
\renewcommand{\headrulewidth}{0.4pt}

\newtheorem{theorem}{Theorem}[section]
\newtheorem{lemma}[theorem]{Lemma}
\newtheorem{corollary}[theorem]{Corollary}
\newtheorem{proposition}[theorem]{Proposition}
\newtheorem{definition}[theorem]{Definition}
\newtheorem{remark}[theorem]{Remark}
\newtheorem{example}[theorem]{Example}
\newtheorem{warning}[theorem]{Warning}

\newcommand{\Q}{\mathbb Q}
\newcommand{\N}{\mathbb N}
\newcommand{\R}{\mathbb R}
\newcommand{\E}{\mathbb E}
\newcommand{\Pp}{\mathbb P}
\newcommand{\Var}{\operatorname{Var}}
\newcommand{\Cov}{\operatorname{Cov}}
\newcommand{\dd}{\,\mathrm d}
\newcommand{\ind}{\mathbf 1}
\newcommand{\PhiN}{\Phi}
\newcommand{\phiN}{\varphi}
\newcommand{\dk}{d_{\mathrm K}}
\newcommand{\abs}[1]{\left|#1\right|}
\newcommand{\norm}[1]{\left\lVert#1\right\rVert}
\newcommand{\ipra}{\mathrm{iPRA}}
\newcommand{\PR}{\mathrm{PR}}
\newcommand{\code}[1]{\nolinkurl{#1}}

\title{\bfseries Constructive Central Limit Theorems\\
       in Intuitionistic Primitive Recursive Arithmetic\\[0.5em]
       \large Berry--Esseen, Lindeberg, Lyapunov, de Moivre--Laplace,\\
       effective extensions, and a Lean formalization route}
\author{Prepared for Vladimir Reshetnikov}
\date{15 November 2026}

\begin{document}
\maketitle

\begin{abstract}
This paper reconstructs a large part of the statement inventory on the
Wikipedia and MathWorld pages about the central limit theorem inside the
finitary, constructive setting of intuitionistic primitive recursive
arithmetic (\(\ipra\)).  The internal objects are finite rational probability
laws, exact finite product spaces, primitive-recursive regular names for real
numbers, and explicit moduli.  No countable probability space, arbitrary
real-valued function, choice principle, or nonconstructive compactness theorem
is used in the core argument.

The main quantitative theorem is a finite triangular-array Berry--Esseen
bound.  If independent centered finite laws have total variance \(s^2>0\) and
\(\rho=\sum_j\E|X_j|^3\), then, at every rational raw threshold \(q\),
\[
 \abs{\Pp\!\left[\sum_j X_j\le q\right]-\PhiN(q/s)}
 \le \min\!\left(1,64\frac{\rho}{s^3}\right).
\]
A separate bounded-increment proof gives the sharper estimate
\(\dk\le\min(1,7B/s)\) when \(|X_j|\le B\).  The proof is elementary and
arithmetically explicit: finite Taylor inequalities, a damping estimate from
an independent copy, a finite product telescoping identity, and Esseen's
smoothing kernel
\(v_T(x)=(1-\cos Tx)/(\pi T x^2)\).  The constants are deliberately
non-optimal but have a complete rational audit.

Consequences include the iid CLT with an explicit \(O(n^{-1/2})\) rate,
uniform CDF convergence, the characteristic-function formulation,
quantitative Lindeberg and Lyapunov theorems, the binomial and local
de Moivre--Laplace theorems, fixed-dimensional multivariate convergence,
an effective random-sum theorem, bounded finite-range dependence, a
finite-tree martingale CLT, lattice local limits, the weak law of large
numbers, and the multiplicative/lognormal CLT.  The unrestricted classical
finite-variance theorem is recovered for effectively presented laws carrying
a primitive-recursive second-moment-tail modulus.

The paper also identifies overbroad source wording: finite variances alone do
not imply a CLT for an arbitrary non-identical array, and a general sample
mean is only asymptotically, not exactly, Gaussian.  A final section gives a
concrete Lean module plan and the exact \(\ipra\) object-language endpoint.
The accompanying Lean files are honest, uncompiled blueprints; they do not
claim a completed kernel-checked derivation.
\end{abstract}

\tableofcontents
\newpage

\section{What is being proved}

\subsection{The constructive replacement for textbook quantifiers}

Classical probability writes ``let \(X\) be an arbitrary real random
variable'' and then quantifies over measures, measurable functions, and
limits.  First-order \(\ipra\) cannot literally quantify over those
higher-type objects.  This is not an obstacle to a constructive central limit
theorem; it determines the correct statement.

A theorem in this paper is a first-order theorem schema for a fixed
primitive-recursive presentation.  The presentation contains finite rational
laws, or a primitive-recursive sequence of finite approximants, together with
all moduli needed to control tails, continuity, and numerical approximation.
The conclusion supplies a primitive-recursive cutoff.  In particular,
``\(F_n\to\PhiN\)'' means
\[
 \forall k\;\forall n\ge N(k)\;\forall q\in\Q,
 \qquad \abs{F_n(q)-\PhiN(q)}\le 2^{-k},                 \tag{1.1}
\]
where \(N\) is a displayed primitive-recursive function and \(\PhiN(q)\) is
a regular interval name.  Internally, the absolute-value inequality is a
family of rational interval comparisons with a dyadic slack.

The use of rational thresholds is essential.  A normalized sum may contain an
irrational square root, whereas the event
\(\sum X_j\le q\) is exactly decidable when the values and \(q\) are rational.
Thus the internally strongest convenient form is
\[
 \abs{\Pp[S\le q]-\PhiN(q/s)}\le\varepsilon
 \qquad(q\in\Q),                                      \tag{1.2}
\]
with \(s=\sqrt{\Var S}\) represented by a regular name.  Externally, because
all jumps of a finite rational law occur at rational points and \(\PhiN\) is
continuous, (1.2) is exactly the usual Kolmogorov-distance assertion.

\subsection{Three levels of coverage}

We distinguish three claims.

\begin{enumerate}
\item \emph{Internal finite theorem.}  All probabilities are exact finite
  sums and the proof directly arithmetizes in \(\ipra\).
\item \emph{Effective completion theorem.}  An infinite or continuous law is
  supplied by finite rational approximants and primitive-recursive tail and
  quantization moduli.  The finite theorem transfers with an explicit error.
\item \emph{Unrestricted classical theorem.}  The usual statement ranges over
  all measures satisfying a qualitative hypothesis.  Such a statement need
  not have a uniform computable modulus and is not, without a presentation,
  a single first-order \(\ipra\) sentence.
\end{enumerate}

The main Berry--Esseen theorem is level 1.  The classical finite-variance iid
CLT is level 2.  The theorem ledger accompanying this paper says which level
is reached for each item on the requested pages.

\subsection{Main results}

Let \(X_1,\ldots,X_m\) be independent finite rational laws, centered so that
\(\E X_j=0\).  Put
\[
 s^2=\sum_{j<m}\E X_j^2>0,
 \qquad \rho=\sum_{j<m}\E|X_j|^3,
 \qquad S=\sum_{j<m}X_j.                               \tag{1.3}
\]

\begin{theorem}[Finite Berry--Esseen theorem]\label{thm:be-main}
For every rational \(q\),
\[
 \abs{\Pp[S\le q]-\PhiN(q/s)}
 \le \min\!\left(1,64\frac{\rho}{s^3}\right).          \tag{1.4}
\]
Every quantity in (1.4) has a uniform primitive-recursive presentation, and
the proof uses only quantifier-free induction after all finite lengths and
precision parameters have been fixed.
\end{theorem}

\begin{theorem}[Bounded array theorem]\label{thm:bounded-main}
If additionally \(|X_j|\le B\) for every support point, then
\[
 \abs{\Pp[S\le q]-\PhiN(q/s)}
 \le \min\!\left(1,7\frac{B}{s}\right)                 \tag{1.5}
\]
for every rational \(q\).
\end{theorem}

There is a typographical convention to note: when a displayed minimum is
used below, it is always harmless to replace the right-hand side by the
second entry if that entry exceeds one.

\begin{corollary}[iid finite-law CLT with rate]\label{cor:iid-rate}
Let \(X\) be a finite rational law with mean \(\mu\), variance
\(\sigma^2>0\), and \(\rho_3=\E|X-\mu|^3\).  For iid copies and
\(S_n=X_1+\cdots+X_n\),
\[
 \sup_x\abs{\Pp\!\left[\frac{S_n-n\mu}{\sigma\sqrt n}\le x\right]
                 -\PhiN(x)}
 \le \min\!\left(1,\frac{64\rho_3}{\sigma^3\sqrt n}\right). \tag{1.6}
\]
The internal version quantifies over rational raw thresholds.  A sufficient
primitive-recursive cutoff for error \(2^{-k}\) is any \(N(k)\) satisfying
\[
 N(k)(\sigma^2)^3\ge 2^{2k+12}\rho_3^2.                \tag{1.7}
\]
\end{corollary}

The factor \(2^{12}=64^2\) makes (1.7) a purely rational test; no division by
an irrational standard deviation occurs in the checker.

\section{Audit of the requested source statements}

The two requested pages~\cite{WikipediaCLT,MathWorldCLT} contain a mixture of precise theorems, informal
interpretations, and claims requiring extra hypotheses.  The following
corrections are mathematically important.

\begin{warning}[Finite variances do not suffice for an arbitrary array]
For each \(n\), take \(X_{n,1}\) to be a Rademacher variable and
\(X_{n,j}=0\) for \(j>1\).  The variables in each row are independent and the
total variance is one, but the normalized sum is Rademacher for every \(n\),
not Gaussian.  A non-identical array therefore needs an infinitesimal,
Lindeberg, Lyapunov, or comparable condition.  The corrected theorem is
proved in \cref{sec:lindeberg}.
\end{warning}

\begin{warning}[A general sample mean is not exactly normal]
For iid non-Gaussian summands, the average has mean \(\mu\) and variance
\(\sigma^2/n\), but its distribution is merely asymptotically Gaussian.
If the summands themselves are Gaussian, characteristic functions show that
the average is exactly \(\mathcal N(\mu,\sigma^2/n)\).  Thus the exact and
asymptotic assertions must not be conflated.
\end{warning}

\begin{warning}[The formal power-series derivation needs repair]
Expanding \(e^{itX}\) into an infinite series and interchanging expectation
with that series requires hypotheses much stronger than finite variance.
Our proof uses the finite identity
\[
 \abs{e^{iu}-1-iu+\tfrac12u^2}\le \frac{|u|^3}{6},     \tag{2.1}
\]
so no infinite interchange is needed.  With only a second moment, truncation
and the Lindeberg argument replace the third-order estimate.
\end{warning}

The page statements proved directly or effectively include the classical
CLT, pointwise and uniform CDF convergence, Berry--Esseen, Lyapunov,
Lindeberg, de Moivre--Laplace, the characteristic-function and convolution
forms, multivariate convergence, random sums, a finite-range dependence
case, a martingale case, the variance formula, products/lognormality, and the
law-of-large-numbers comparison.  The high-dimensional convex-body and
random-geometry theorems use separate machinery and are explicitly listed as
out of scope rather than being silently claimed.

\section{Coding finite probability in \texorpdfstring{\(\ipra\)}{iPRA}}\label{sec:coding}

\subsection{Finite rational laws}

\begin{definition}[Law code]
A finite rational law is a tuple
\[
 L=(r;(a_j,b_j,w_j)_{j<r};W)                           \tag{3.1}
\]
with \(b_j>0\), \(w_j\in\N\), \(W=\sum_{j<r}w_j>0\),
and value \(x_j=a_j/b_j\).  Repetitions are allowed.  For a rational-valued
observable \(f\),
\[
 \E_L f=\frac1W\sum_{j<r}w_j f(x_j).                  \tag{3.2}
\]
\end{definition}

Signed integers and positive-denominator rationals are coded by natural
pairs.  Addition, multiplication, absolute value, comparison, finite sums,
and bounded maxima are primitive recursive.  The validity predicate for
(3.1) is decidable.

A product of laws \(L_0,\ldots,L_{m-1}\) is enumerated by a mixed-radix
index.  The weight of an outcome tuple is the product of its coordinate
weights and the total weight is the product of the coordinate totals.
Consequently independence is not an axiom: it is an exact identity of finite
sums.

\begin{lemma}[Finite independence algebra]\label{lem:finite-algebra}
For observables depending on disjoint coordinates,
\[
 \E[f(X_A)g(X_B)]=\E f(X_A)\,\E g(X_B).                \tag{3.3}
\]
Hence centered independent variables satisfy
\[
 \E\sum_jX_j=0,
 \qquad \Var\!\left(\sum_jX_j\right)=\sum_j\Var(X_j), \tag{3.4}
\]
and their characteristic functions multiply.
\end{lemma}

\begin{proof}
Expand both sides of (3.3) over the mixed-radix product.  The numerator is
\[
 \sum_{u,v}w_A(u)w_B(v)f(u)g(v)
 =\left(\sum_uw_A(u)f(u)\right)
  \left(\sum_vw_B(v)g(v)\right),                      \tag{3.5}
\]
and the denominator factors in the same way.  Formula (3.4) follows by
expanding the square; all cross terms vanish by (3.3) and centering.  The
characteristic-function identity is the same calculation with finite
regular names for sine and cosine.  Each equality is a bounded fold and is
proved by quantifier-free induction on a fixed finite length.
\end{proof}

\subsection{Regular interval names}

\begin{definition}[Regular name]
A real number \(x\) is presented by primitive-recursive rational functions
\(L_x(p),U_x(p)\) such that
\[
 L_x(p)\le L_x(p+1)\le U_x(p+1)\le U_x(p),
 \qquad U_x(p)-L_x(p)\le2^{-p}.                        \tag{3.6}
\]
\end{definition}

Equality and order are never decided by an oracle.  A theorem outputs
interval inclusions with an explicit dyadic slack.  A positive real used as
a denominator is accompanied by a rational margin \(c>0\) and a proof
\(c\le x\).  Addition and bounded multiplication are interval arithmetic;
reciprocal and square root use rational bisection.  Complex numbers are pairs
of regular names.

This is the same foundational convention used by the constructive analysis
packets in the ProveIt repository.  It is important that the program is
fixed.  \(\ipra\) can reason about every fixed primitive-recursive evaluator,
but it cannot contain a universal evaluator for all primitive-recursive
program codes without a diagonal contradiction.

\subsection{Primitive-recursive convergence}

For a presented sequence \(x_n\) and limit \(x\), the assertion
\(x_n\to x\) is always accompanied by a primitive-recursive modulus
\(N(k)\):
\[
 n\ge N(k)\quad\Longrightarrow\quad |x_n-x|\le2^{-k}. \tag{3.7}
\]
For distribution functions, the threshold is a separate number variable and
(3.7) is uniform in its code.  Qualitative phrases such as ``choose \(n\)
large enough'' are eliminated from the internal theorem.

\section{Constructive Gaussian and Fourier toolkit}\label{sec:gaussian}

\subsection{Elementary analytic names}

The required constants and functions have standard finite algorithms.

\begin{itemize}
\item \(\pi\) is named by a Machin formula, for example
\[
 \pi=16\arctan(1/5)-4\arctan(1/239),                  \tag{4.1}
\]
with alternating-series tail bounds.
\item \(\exp,\sin,\cos\) use range reduction and finite Taylor sums with
factorial remainders.  The Taylor/Cauchy packet supplies the general finite
certificate.
\item Positive square roots use bisection in a rational bracket.
\item Integrals on compact rational intervals are named by finite Riemann
sums with an explicit derivative or Lipschitz bound.  Improper Gaussian
integrals are a compact integral plus a Mills-type tail estimate.
\end{itemize}

Only two coarse rational bounds on \(\pi\) enter the constant audit:
\[
 \frac{25}{8}<\pi<\frac{16}{5}.                       \tag{4.2}
\]
Both follow from short finite arctangent sums.

\subsection{The normal density and CDF}

Define
\[
 \phiN(x)=\frac{1}{\sqrt{2\pi}}e^{-x^2/2},
 \qquad
 \PhiN(x)=\frac12+\int_0^x\phiN(u)\dd u.              \tag{4.3}
\]
The Gaussian integral is obtained constructively from the Wallis/Stirling
machinery, so \(\int_{\R}\phiN=1\).  Since the maximum occurs at zero,
(4.2) gives
\[
 0\le\phiN(x)\le\frac1{\sqrt{2\pi}}<\frac25.          \tag{4.4}
\]
Thus \(\PhiN\) has the global rational Lipschitz constant \(2/5\).
For \(x>0\), integration by parts gives the effective Mills bound
\[
 1-\PhiN(x)\le\frac{\phiN(x)}{x},                     \tag{4.5}
\]
and symmetry treats the left tail.

\begin{lemma}[Gaussian characteristic function]\label{lem:gauss-cf}
The characteristic function of (4.3) is
\[
 \int_{\R}e^{itx}\phiN(x)\dd x=e^{-t^2/2}.            \tag{4.6}
\]
\end{lemma}

\begin{proof}
The sine integral vanishes by oddness.  Let
\(g(t)=\int\cos(tx)\phiN(x)\dd x\).  Truncate at \([-A,A]\), differentiate
the finite integral, and integrate by parts using
\(\phiN'(x)=-x\phiN(x)\).  Boundary terms are bounded by
\(2\phiN(A)\), so the explicit Gaussian tail lets \(A\) be chosen
primitive-recursively.  One obtains \(g'(t)=-tg(t)\) and \(g(0)=1\).
Multiplying by \(e^{t^2/2}\) gives a function with derivative zero; the
finite mean-value estimate proves it is constantly one.  Every limiting step
has the displayed tail modulus, so this is a regular-name argument, not an
appeal to abstract completeness.
\end{proof}

\subsection{Three finite inequalities}

The characteristic-function proof uses only the following rationally
certified inequalities:
\begin{align}
 \abs{e^{iu}-1-iu+\tfrac12u^2}&\le |u|^3/6,           \tag{4.7}\\
 0\le e^{-x}-1+x&\le x^2/2 \quad(x\ge0),             \tag{4.8}\\
 1-\cos u&\ge u^2/3 \quad(|u|\le1).                  \tag{4.9}
\end{align}
The first two are Taylor formulas with integral or alternating remainder.
For (4.9), the alternating cosine series gives
\(1-\cos u\ge u^2/2-u^4/24\ge11u^2/24>u^2/3\).
All inequalities can be cleared of denominators and checked with QF
induction on the finite Taylor order.

\section{Esseen smoothing with an explicit kernel}\label{sec:smoothing}

The passage from characteristic functions to distribution functions is the
only place where discontinuous indicators enter.  We give the full argument
because its finitary form is especially suitable for \(\ipra\).

For \(T>0\), define the even kernel
\[
 v_T(x)=
 \begin{cases}
 \displaystyle\frac{1-\cos(Tx)}{\pi T x^2},&x\ne0,\\[0.7em]
 \displaystyle\frac{T}{2\pi},&x=0.
 \end{cases}                                           \tag{5.1}
\]
Its value at zero is the removable limit.  The elementary Dirichlet-integral
calculation gives
\[
 v_T(x)\ge0,\qquad \int_{\R}v_T(x)\dd x=1,             \tag{5.2}
\]
and its Fourier transform is triangular:
\[
 \widehat v_T(t)=
 \begin{cases}
 1-|t|/T,&|t|\le T,\\
 0,&|t|>T.
 \end{cases}                                           \tag{5.3}
\]
Moreover, because \(1-\cos(Tx)\le2\),
\[
 \int_{|x|>h}v_T(x)\dd x\le\frac4{\pi T h}.           \tag{5.4}
\]

For constructivity, (5.2)--(5.3) are not black boxes.  Here is an
explicit modulus for the only conditionally convergent integral.  For
\(a>0\), put
\[
 I(a)=\int_0^\infty e^{-au}\frac{\sin u}{u}\dd u.
\]
Differentiation is first performed on \([0,A]\); the exponential tail is
uniform, so the limit is legitimate with a displayed bound.  Since
\[
 I'(a)=-\int_0^\infty e^{-au}\sin u\dd u
       =-\frac1{1+a^2},
 \qquad \lim_{a\to\infty}I(a)=0,
\]
one obtains \(I(a)=\arctan(1/a)\).  On the other hand, Dirichlet integration
by parts gives, for \(B>A>0\),
\[
 \left|\int_A^B\frac{\sin u}{u}\dd u\right|\le\frac2A,
 \qquad
 \left|\int_A^\infty e^{-au}\frac{\sin u}{u}\dd u\right|
 \le\frac2A.
\]
On \([0,A]\), the damped and undamped integrals differ by at most \(aA\),
because \(1-e^{-au}\le au\).  Thus choosing first rational \(A\) and then
rational \(a\) makes the total comparison error at most
\(4/A+aA\).  Since \(\arctan(1/a)\to\pi/2\) with its alternating-series
modulus, this proves
\(\int_0^\infty\sin u/u\dd u=\pi/2\) constructively.

An integration by parts, whose boundary term is bounded by \(2/x\) at
infinity and by \(T^2x/2\) at zero, now gives
\[
 \int_0^\infty\frac{1-\cos(Tx)}{x^2}\dd x
 =T\int_0^\infty\frac{\sin(Tx)}x\dd x
 =\frac{\pi T}{2}.
\]
This proves the normalization in (5.2).  More importantly, direct finite
integration yields the inverse-transform identity
\[
 v_T(x)=\frac1{2\pi}\int_{-T}^{T}
       \left(1-\frac{|t|}{T}\right)e^{itx}\dd t.     \tag{5.5}
\]
For \(x\ne0\), the real part is
\(\pi^{-1}\int_0^T(1-t/T)\cos(tx)\dd t
 =(1-\cos(Tx))/(\pi T x^2)\), and the imaginary part vanishes; the value at
zero is the area of the triangle divided by \(2\pi\).  Formula (5.3) is the
corresponding finite Fourier-pair statement.  In the arithmetic proof one
uses (5.5) directly, with finite Riemann enclosures, so no general Fourier
inversion theorem is assumed.  Thus every integral here is a compact finite
sum plus one of the explicit remainders just displayed.

\begin{theorem}[Constructive Esseen smoothing]\label{thm:esseen}
Let \(F\) and \(G\) be distribution functions with finite first moments,
characteristic functions \(f\) and \(g\), and suppose that \(G\) is
Lipschitz with constant \(M\).  For every \(T>0\),
\[
 \sup_x|F(x)-G(x)|
 \le \frac1\pi\int_{-T}^{T}
        \abs{\frac{f(t)-g(t)}{t}}\dd t
      +\frac{24M}{\pi T}.                              \tag{5.6}
\]
For finite rational \(F\), the supremum may be replaced internally by the
universal assertion over rational \(x\).
\end{theorem}

\begin{proof}
Put \(\Delta=F-G\), \(\eta=\sup_x|\Delta(x)|\), and
\(\Delta_T=\Delta*v_T\), \(\eta_T=\sup_x|\Delta_T(x)|\).
We first show
\[
 \eta\le2\eta_T+\frac{24M}{\pi T}.                    \tag{5.7}
\]
For a finite rational law, \(\Delta\) is a continuous function between
finitely many jumps and tends to zero at both ends, so a one-sided extremum
is attained.  For a general effective approximant, replace the maximum by a
point within a supplied rational \(\epsilon\) and carry the \(\epsilon\)
through; the final slack absorbs it.

Suppose first that \(\Delta(x_0)=\eta>0\).  Monotonicity of \(F\) and the
Lipschitz bound for \(G\) give, for \(s>0\),
\[
 \Delta(x_0+s)\ge\eta-Ms.                              \tag{5.8}
\]
Set \(h=\eta/(2M)\) and \(y=x_0+h\).  If \(|x|\le h\), then with
\(s=h-x\),
\[
 \Delta(y-x)\ge\eta/2+Mx.                             \tag{5.9}
\]
Outside this interval we only use \(\Delta\ge-\eta\).  Since \(v_T\) is
even, the integral of \(Mxv_T(x)\) over \([-h,h]\) vanishes.  Applying
(5.4),
\begin{align*}
 \eta_T&\ge\Delta_T(y)\\
 &\ge \frac\eta2\left(1-\frac4{\pi Th}\right)
       -\eta\frac4{\pi Th}
  =\frac\eta2-\frac{12M}{\pi T}.                     \tag{5.10}
\end{align*}
This is (5.7).  If the negative extremum has larger magnitude, reflect both
laws by \(x\mapsto-x\), which preserves the hypotheses and the convolution
norm.

It remains to estimate \(\eta_T\).  Fourier inversion for the smoothed CDF
difference, using (5.3), gives
\[
 \Delta_T(x)=\frac{1}{2\pi}\int_{-T}^{T}
  \frac{-e^{-itx}}{it}\,[f(t)-g(t)]
  \left(1-\frac{|t|}{T}\right)\dd t.                 \tag{5.11}
\]
The apparent singularity at zero is removable because finite first moments
give \(f(t)-g(t)=O(t)\).  Hence
\[
 \eta_T\le\frac1{2\pi}\int_{-T}^{T}
       \abs{\frac{f(t)-g(t)}{t}}\dd t.                \tag{5.12}
\]
Substitution into (5.7) proves (5.6).

For \(\ipra\), (5.11) is implemented on a finite rational mesh.  The
Lipschitz moduli of the finite characteristic functions and of the kernel
supply a Riemann-sum error; the tails have already disappeared because of
(5.3).  The proof therefore reduces to finite rational inequalities and does
not use a measure-theoretic inversion theorem as an oracle.
\end{proof}

\section{A bounded finite-array theorem}\label{sec:bounded}

We first prove the sharper bounded result.  Besides being useful in its own
right, it is the engine for the constructive Lindeberg theorem.

Normalize so that the total variance is one.  Write
\[
 v_j=\E Y_j^2,\qquad \beta_j=\E|Y_j|^3,
 \qquad \sum_jv_j=1,                                  \tag{6.1}
\]
where the \(Y_j\) are independent, centered, and \(|Y_j|\le b\).  Let
\(a_j(t)=\E e^{itY_j}\) and \(g_j(t)=e^{-v_jt^2/2}\).

\begin{lemma}[Damping from an independent copy]\label{lem:bounded-damping}
If \(|t|\le1/(2b)\), then
\[
 |a_j(t)|\le e^{-v_jt^2/3}.                            \tag{6.2}
\]
\end{lemma}

\begin{proof}
Let \(Y'_j\) be an independent copy.  Finite independence gives
\[
 |a_j(t)|^2=\E e^{it(Y_j-Y'_j)}
            =\E\cos(t(Y_j-Y'_j)).                     \tag{6.3}
\]
Because \(|t(Y_j-Y'_j)|\le1\), inequality (4.9) yields
\[
 1-|a_j(t)|^2
 \ge\frac{t^2}{3}\E(Y_j-Y'_j)^2
 =\frac{2v_jt^2}{3}.                                  \tag{6.4}
\]
Now \(1-u\le e^{-u}\) and square roots give (6.2).
\end{proof}

\begin{lemma}[One-factor replacement]\label{lem:one-factor}
For every real \(t\),
\[
 |a_j(t)-g_j(t)|
 \le\frac{\beta_j|t|^3}{6}+\frac{v_j^2t^4}{8}.        \tag{6.5}
\]
\end{lemma}

\begin{proof}
Taking expectations in (4.7), using \(\E Y_j=0\), gives
\[
 \abs{a_j(t)-1+\frac12v_jt^2}
 \le\frac16\beta_j|t|^3.                             \tag{6.6}
\]
Apply (4.8) with \(x=v_jt^2/2\) and add the two errors.
\end{proof}

\begin{theorem}[Bounded array, normalized form]\label{thm:bounded-normalized}
If \(\sum_j\E Y_j^2=1\) and \(|Y_j|\le b\), then
\[
 \dk\!\left(\sum_jY_j,\mathcal N(0,1)\right)
 \le\min(1,7b).                                       \tag{6.7}
\]
\end{theorem}

\begin{proof}
If \(b\ge1/7\), the trivial bound \(\dk\le1\) is enough.  Assume
\(b<1/7\) and set \(T=1/(2b)\).  The finite product identity
\[
 \prod_ja_j-\prod_jg_j
 =\sum_j(a_j-g_j)\prod_{k<j}a_k\prod_{k>j}g_k         \tag{6.8}
\]
is proved by induction on the number of factors.  Both (6.2) and
\(g_k\le e^{-v_kt^2/3}\) hold on \([-T,T]\).  Since
\(v_j\le b^2<1/49\), the product excluding the \(j\)-th factor is at most
\[
 e^{-(1-v_j)t^2/3}\le e^{-t^2/4}.                     \tag{6.9}
\]
Also \(\sum_j\beta_j\le b\sum_jv_j=b\), and
\(\sum_jv_j^2\le b^2\sum_jv_j=b^2\).  Therefore (6.5)--(6.9) imply
\[
 \abs{\E e^{it\sum_jY_j}-e^{-t^2/2}}
 \le e^{-t^2/4}
       \left(\frac{b|t|^3}{6}+\frac{b^2t^4}{8}\right)
 \le\frac{11}{48}b|t|^3e^{-t^2/4}.                   \tag{6.10}
\]
The last inequality uses \(b|t|\le1/2\).

Apply \cref{thm:esseen} with the standard normal, whose density maximum
\(M\) is less than \(2/5\).  Since
\[
 \int_{-\infty}^{\infty}t^2e^{-t^2/4}\dd t=4\sqrt\pi, \tag{6.11}
\]
we obtain
\[
 \dk\le \frac{11b}{12\sqrt\pi}+\frac{48Mb}{\pi}.     \tag{6.12}
\]
The rational estimates \(1/\sqrt\pi<3/5\), \(M<2/5\), and
\(1/\pi<8/25\) give
\[
 \frac{11}{20}b+\frac{768}{125}b
 =\frac{3347}{500}b<7b.                               \tag{6.13}
\]
All constants are rational consequences of (4.2).
\end{proof}

Rescaling \(Y_j=X_j/s\) proves \cref{thm:bounded-main}.  Notice that the
event is still evaluated as \(\sum X_j\le q\); only the Gaussian argument
uses the regular name \(q/s\).

\section{Finite Berry--Esseen theorem}\label{sec:berry-esseen}

The bounded theorem depends on the largest support value.  A true
Berry--Esseen bound depends only on the sum of third absolute moments.  The
same finite Fourier method gives a clean, deliberately coarse constant 64.

Again normalize \(\sum_jv_j=1\), and put
\[
 \beta_j=\E|Y_j|^3,\qquad \beta=\sum_j\beta_j.         \tag{7.1}
\]
By the finite Lyapunov inequality,
\[
 v_j^{3/2}\le\beta_j.                                  \tag{7.2}
\]

\begin{lemma}[Truncated damping]\label{lem:truncated-damping}
For \(t\in\R\), define
\[
 q_j(t)=\bigl(2v_j-8|t|\beta_j\bigr)_+.               \tag{7.3}
\]
Then
\[
 |a_j(t)|\le e^{-q_j(t)t^2/6}.                         \tag{7.4}
\]
\end{lemma}

\begin{proof}
With \(Z=Y_j-Y'_j\), formula (6.3) remains valid.  On \(|tZ|\le1\), use
(4.9):
\[
 1-|a_j(t)|^2
 \ge\frac{t^2}{3}\E[Z^2\ind_{\{|tZ|\le1\}}].        \tag{7.5}
\]
The discarded second moment satisfies
\[
 \E[Z^2\ind_{\{|tZ|>1\}}]
 \le |t|\E|Z|^3
 \le8|t|\beta_j,                                      \tag{7.6}
\]
because \(|u-v|^3\le4(|u|^3+|v|^3)\).  Since
\(\E Z^2=2v_j\), (7.5)--(7.6) give
\(1-|a_j|^2\ge t^2q_j/3\).  Apply \(1-u\le e^{-u}\)
and take square roots.
\end{proof}

\begin{theorem}[Berry--Esseen, normalized form]\label{thm:be-normalized}
For independent centered finite laws with total variance one,
\[
 \dk\!\left(\sum_jY_j,\mathcal N(0,1)\right)
 \le\min(1,64\beta).                                  \tag{7.7}
\]
\end{theorem}

\begin{proof}
If \(\beta\ge1/64\), the right-hand side is at least one.  Assume
\(0<\beta<1/64\) and set \(T=1/(16\beta)\).  For \(|t|\le T\),
\[
 \sum_kq_k(t)\ge2-8\beta|t|\ge\frac32.                \tag{7.8}
\]
Moreover, (7.2) gives
\(v_j\le\beta_j^{2/3}\le\beta^{2/3}<1/16\), and hence
\[
 \sum_{k\ne j}q_k(t)\ge\frac32-2v_j\ge\frac{11}{8}. \tag{7.9}
\]
The Gaussian factor \(g_k=e^{-v_kt^2/2}\) is also bounded by
\(e^{-q_kt^2/6}\), since \(q_k\le2v_k\).  Therefore every product
excluding the \(j\)-th factor in (6.8) is at most
\[
 e^{-11t^2/48}\le e^{-t^2/5}.                         \tag{7.10}
\]
Using \cref{lem:one-factor}, together with
\[
 \sum_jv_j^2\le\sum_jv_j^{3/2}\le\sum_j\beta_j=\beta,\tag{7.11}
\]
we get
\[
 \abs{\E e^{it\sum_jY_j}-e^{-t^2/2}}
 \le\beta\left(\frac{|t|^3}{6}+\frac{t^4}{8}\right)e^{-t^2/5}
 \qquad(|t|\le T).                                    \tag{7.12}
\]

The integral term in \cref{thm:esseen} is at most
\[
 \frac{2\beta}{\pi}
 \left[
  \frac16\int_0^\infty t^2e^{-t^2/5}\dd t
 +\frac18\int_0^\infty t^3e^{-t^2/5}\dd t
 \right].                                              \tag{7.13}
\]
Integration by parts and the Gaussian integral give
\[
 \int_0^\infty t^2e^{-t^2/5}\dd t=\frac{5\sqrt{5\pi}}4,
 \qquad
 \int_0^\infty t^3e^{-t^2/5}\dd t=\frac{25}{2}.      \tag{7.14}
\]
Using \(\pi>3\) and \(5\pi<16\), (7.13) is less than
\(8\beta/5\).  The smoothing remainder is
\[
 \frac{24M}{\pi T}=\frac{384M}{\pi}\beta
 <\frac{6144}{125}\beta.                              \tag{7.15}
\]
Thus
\[
 \dk<\left(\frac85+\frac{6144}{125}\right)\beta
 =\frac{6344}{125}\beta<64\beta.                     \tag{7.16}
\]
This proves the theorem.
\end{proof}

Rescaling immediately proves \cref{thm:be-main}.  The proof never takes a
logarithm of a complex characteristic function; this removes a branch-choice
issue and makes the arithmetic reconstruction substantially simpler.

\section{The iid theorem, sample means, and effective finite variance}
\label{sec:iid}

\subsection{The finite-law theorem and its cutoff}

\begin{proof}[Proof of \cref{cor:iid-rate}]
Apply \cref{thm:be-main} to
\(X_j-\mu\), \(0\le j<n\).  Then
\[
 s^2=n\sigma^2,
 \qquad
 \rho=n\rho_3,
 \qquad
 64\frac{\rho}{s^3}
 =\frac{64\rho_3}{\sigma^3\sqrt n}.                  \tag{8.1}
\]
This is (1.6).  If (1.7) holds, then after squaring the nonnegative
quantities one gets
\[
 \frac{64^2\rho_3^2}{n\sigma^6}\le2^{-2k},           \tag{8.2}
\]
so the error is at most \(2^{-k}\).  The same conclusion for the sample
mean follows from the identity
\[
 \frac{S_n-n\mu}{\sigma\sqrt n}
 =\frac{\sqrt n}{\sigma}(\overline X_n-\mu).          \tag{8.3}
\]
\end{proof}

Because (8.1) bounds every threshold simultaneously, it proves both the
pointwise CDF statement and the uniform CDF statement on the requested
pages.  No completeness argument turning pointwise convergence into uniform
convergence is needed.  It also separates two assertions that are often
blurred:
\[
 \E\overline X_n=\mu,
 \qquad
 \Var(\overline X_n)=\frac{\sigma^2}{n}               \tag{8.4}
\]
are exact finite identities, whereas normality of \(\overline X_n\) is in
general only asymptotic.  If \(X\) is Gaussian, its characteristic function
is \(e^{i\mu t-\sigma^2t^2/2}\); multiplication of \(n\) such factors proves
that (8.4) is then the mean and variance of an \emph{exact} Gaussian law.

\subsection{A second-moment-tail presentation}

The classical Lindeberg--L\'evy theorem assumes only a finite second moment.
A constructive version must say how that moment is known to be finite.
The following input format is sufficient and covers discrete and continuous
computable laws.

\begin{definition}[Effective finite-variance law]
An effective finite-variance presentation consists of a centered presented
law \(X\), a positive regular name \(\sigma^2\), and, for every \(r\), a
finite rational law \(X^{[r]}\) together with a primitive-recursive coupling
presentation satisfying
\[
 \E\bigl[(X-X^{[r]})^2\bigr]\le \tau(r),
 \qquad \tau(r)\le2^{-r},                             \tag{8.5}
\]
for a primitive-recursive function \(\tau\).  The coupling presentation
supplies rational interval enclosures for its finite approximants and a
modulus making (8.5) an arithmetic certificate.  The means, variances, and
third moments of \(X^{[r]}\) are computed exactly.  Truncation followed by
rational quantization is the basic example.
\end{definition}

The coupling is input data, not an invocation of an abstract optimal-coupling
theorem.  At every requested precision, its probability and moment estimates
are reduced to a finite rational approximation before the iPRA derivation is
constructed.

\begin{lemma}[Transfer between close sums]\label{lem:l2-transfer}
Let \(X_i,Y_i\) be independent pairs with identical pair law, and center both
coordinates.  If \(\E(X-Y)^2\le\delta^2\), then, for every \(h>0\),
\[
 \dk\!\left(\frac{\sum_{i<n}X_i}{\sigma\sqrt n},N(0,1)\right)
 \le
 \dk\!\left(\frac{\sum_{i<n}Y_i}{\sigma\sqrt n},N(0,1)\right)
 +\frac{\delta^2}{\sigma^2h^2}+\frac{2h}{5}.          \tag{8.6}
\]
\end{lemma}

\begin{proof}
Put \(R_n=\sum(X_i-Y_i)\).  Independence and centering give
\(\E R_n^2\le n\delta^2\).  Chebyshev's finite-sum proof yields
\[
 \Pp[|R_n|>h\sigma\sqrt n]\le\frac{\delta^2}{\sigma^2h^2}. \tag{8.7}
\]
On the complementary event, replacing the \(Y\)-sum by the \(X\)-sum moves
a normalized threshold by at most \(h\).  Apply the two inclusions
\(\{Y\le x-h\}\subseteq\{X\le x\}\subseteq\{Y\le x+h\}\), up to the bad
event.  The normal CDF moves by at most \(2h/5\), by (4.4).  This proves the stated transfer estimate.  In a level-2 presentation,
every probability in this argument is enclosed by the supplied finite
coupling approximants with their explicit modulus.
\end{proof}

\begin{theorem}[Effective Lindeberg--L\'evy CLT]\label{thm:effective-iid}
Every effective finite-variance law with \(\sigma^2>0\) has a
primitive-recursive modulus for
\[
 \frac{X_1+\cdots+X_n}{\sigma\sqrt n}
 \Longrightarrow\mathcal N(0,1).                     \tag{8.8}
\]
The same holds after restoring a nonzero mean.
\end{theorem}

\begin{proof}
Fix an error request \(2^{-k}\).  First choose \(r\) so that the right side
of (8.6), with \(\delta^2=4\tau(r)\) and a dyadic \(h\), is below
\(2^{-k-2}\).  The factor four allows for centering both approximants.
The variance \(\sigma_r^2\) of \(X^{[r]}\) then differs from \(\sigma^2\) by
an explicitly controlled amount, because
\[
 |\norm{X}_2-\norm{X^{[r]}}_2|
 \le\norm{X-X^{[r]}}_2.                              \tag{8.9}
\]
Choose \(r\) larger if necessary so that rescaling from \(\sigma_r\) to
\(\sigma\) changes the normal CDF by at most \(2^{-k-2}\).  Finally apply
\cref{cor:iid-rate} to the fixed finite law \(X^{[r]}\), using its exactly
computed third moment, and choose \(n\) from (1.7) so that its Berry--Esseen
error is below \(2^{-k-2}\).  The triangle inequality gives (8.8).

All choices are bounded searches through rational interval certificates.
There is no extraction of a modulus from the bare assertion
\(\E X^2<\infty\): the tail/quantization modulus is part of the constructive
hypothesis.
\end{proof}

\subsection{Characteristic functions and convolution}

For a centered finite law of variance one, (4.7) with a truncation argument
at order two gives, for every fixed rational \(t\),
\[
 \phi_X(t/\sqrt n)
 =1-\frac{t^2}{2n}+r_n(t),
 \qquad n|r_n(t)|\longrightarrow0                   \tag{8.10}
\]
with a primitive-recursive modulus supplied by the second-moment-tail data.
The finite product identity then yields
\[
 \phi_{(X_1+\cdots+X_n)/\sqrt n}(t)
 =\phi_X(t/\sqrt n)^n\longrightarrow e^{-t^2/2}.     \tag{8.11}
\]
For finite laws with a third moment, (7.12) is already a quantitative and
uniform version on a growing interval.  Esseen smoothing proves the CDF
conclusion directly, so an abstract L\'evy continuity theorem is unnecessary.

At the distribution level, independence says that the law of a sum is the
finite convolution of the laws.  Hence (8.8) is equivalently a theorem about
normalized convolution powers.  This proves the characteristic-function and
convolution formulations on the source pages under the exact constructive
hypotheses stated here.

\section{Lindeberg and Lyapunov central limit theorems}
\label{sec:lindeberg}

Let a row consist of independent centered finite laws \(Y_0,\ldots,Y_{m-1}\)
with total variance one.  Define the normalized Lindeberg tail
\[
 L(\eta)=\sum_{j<m}\E\!
   \left[Y_j^2\ind_{\{|Y_j|>\eta\}}\right].           \tag{9.1}
\]

\begin{theorem}[Quantitative finite Lindeberg theorem]
\label{thm:quant-lindeberg}
Suppose \(0<\lambda\le1/2\), \(\eta>0\), and
\(L(\eta)\le\lambda^3\).  Then
\[
 \dk\!\left(\sum_jY_j,\mathcal N(0,1)\right)
 \le 28\eta+3\lambda.                                \tag{9.2}
\]
\end{theorem}

\begin{proof}
Truncate and recenter:
\[
 T_j=Y_j\ind_{\{|Y_j|\le\eta\}}
       -\E\!\bigl[Y_j\ind_{\{|Y_j|\le\eta\}}\bigr],
 \qquad R_j=Y_j-T_j.                                  \tag{9.3}
\]
Then \(|T_j|\le2\eta\), both families are centered, and independence gives
\[
 \E\left(\sum_jR_j\right)^2
 =\sum_j\Var(R_j)\le L(\eta).                        \tag{9.4}
\]
Let \(q=\Var(\sum_jT_j)\).  The triangle inequality in the finite
\(L^2\) norm gives
\[
 |\sqrt q-1|\le\sqrt{L(\eta)}\le\lambda,             \tag{9.5}
\]
where \(\lambda^{3/2}\le\lambda\).  Thus \(\sqrt q\ge1/2\) and
\(|T_j|/\sqrt q\le4\eta\).  By \cref{thm:bounded-normalized},
\[
 \dk\!\left(\frac{\sum_jT_j}{\sqrt q},N(0,1)\right)
 \le28\eta.                                          \tag{9.6}
\]
The two normal CDFs \(\PhiN(x/\sqrt q)\) and \(\PhiN(x)\) differ uniformly
by at most \(\lambda/2\).  Indeed, the derivative with respect to the scale
is \(|y|\phiN(y)/a\), and for \(a\ge1/2\) the elementary bound
\(\sup_y|y|\phiN(y)<1/4\) gives \(1/2\).

Put \(R=\sum_jR_j\).  In the two one-sided comparisons between \(T+R\) and
\(T\), move the threshold by \(\lambda\).  Chebyshev and (9.4) give
\(\Pp(|R|>\lambda)\le\lambda\), while the Gaussian shift costs less than
\(2\lambda/5\).  Together with (9.6) and the scale cost this is less than
\(28\eta+2\lambda\), and hence proves the slightly looser displayed bound.
\end{proof}

\begin{corollary}[Effective Lindeberg--Feller theorem]
\label{cor:lindeberg}
For a primitive-recursive triangular array with total row variance one,
assume that a primitive-recursive function supplies, for every dyadic
\(\eta>0\) and precision \(k\), a row cutoff after which
\(L_n(\eta)\le2^{-k}\).  Then the row sums converge uniformly in CDF to the
standard normal law, with a primitive-recursive modulus.
\end{corollary}

\begin{proof}
Given \(2^{-k}\), choose dyadic \(\eta\) with \(28\eta<2^{-k-1}\), then
dyadic \(\lambda\le1/2\) with \(3\lambda<2^{-k-1}\).  Use the supplied
Lindeberg modulus until \(L_n(\eta)\le\lambda^3\), and apply
\cref{thm:quant-lindeberg}.
\end{proof}

Lindeberg also implies the Feller infinitesimal condition.  Indeed, for each
summand and every \(\eta>0\),
\[
 \E Y_j^2\le\eta^2+L(\eta),                           \tag{9.7}
\]
so the largest individual variance tends to zero when the row Lindeberg
quantities do.

\subsection{Lyapunov's condition}

For a fixed positive rational \(\delta\), put
\[
 \Lambda_\delta=\sum_j\E|Y_j|^{2+\delta}.             \tag{9.8}
\]
Rational powers are regular names obtained by bounded root bisection; for
integer \(\delta\) everything is rational outright.  On \(|Y_j|>\eta\),
\(Y_j^2\le |Y_j|^{2+\delta}/\eta^\delta\), hence
\[
 L(\eta)\le\eta^{-\delta}\Lambda_\delta.             \tag{9.9}
\]
Taking
\(\eta=\Lambda_\delta^{1/(\delta+3)}\) and
\(\lambda=(\Lambda_\delta/\eta^\delta)^{1/3}=\eta\)
in (9.2), whenever \(\eta\le1/2\), gives
\[
 \dk\le31\Lambda_\delta^{1/(\delta+3)}.             \tag{9.10}
\]
Thus a primitive-recursive Lyapunov modulus implies a CLT modulus.  For the
most common case \(\delta=1\), \cref{thm:be-normalized} gives the much
stronger estimate
\[
 \dk\le\min(1,64\Lambda_1).                           \tag{9.11}
\]
This proves both the Lyapunov theorem and the implication
``Lyapunov \(\Rightarrow\) Lindeberg'' stated on the source page.

\subsection{The converse fails constructively}

\begin{example}[Lindeberg without Lyapunov]\label{ex:nonconverse}
For row \(m\ge2\), take \(m^2-m\) independent Rademacher variables of
amplitude \(1/m\), and one independent spike variable \(Z_m\) with
\[
 \Pp(Z_m=m^2)=\Pp(Z_m=-m^2)=\frac1{2m^5},
 \qquad \Pp(Z_m=0)=1-\frac1{m^5}.                    \tag{9.12}
\]
The row variance is
\[
 (m^2-m)m^{-2}+m^4m^{-5}=1.                           \tag{9.13}
\]
For every fixed \(\eta>0\), the small variables eventually contribute no
Lindeberg tail and the spike contributes \(1/m\), so (9.1) tends to zero.
But the row third absolute moment contains
\(\E|Z_m|^3=m\), so the \(\delta=1\) Lyapunov quantity diverges.  All
probabilities are rational and all verifications are finite identities in
\(\ipra\).
\end{example}

\section{The de Moivre--Laplace theorems}
\label{sec:demoivre}

Let \(B_n\sim\operatorname{Bin}(n,p)\), where
\(0<p<1\) is rational, and put \(q=1-p\).  Bernoulli's third centered
absolute moment is
\[
 \rho_3=pq(p^2+q^2),
 \qquad \sigma^2=pq.                                  \tag{10.1}
\]
Thus \cref{cor:iid-rate} immediately gives the global theorem.

\begin{theorem}[Global de Moivre--Laplace]
For every integer \(k\),
\[
 \left|
 \Pp(B_n\le k)-
 \PhiN\!\left(\frac{k-np}{\sqrt{npq}}\right)
 \right|
 \le\min\!\left(1,
       \frac{64(p^2+q^2)}{\sqrt{npq}}\right).         \tag{10.2}
\]
With the continuity correction \(k+1/2\), the right side may be increased by
\(1/(5\sqrt{npq})\).
\end{theorem}

\begin{proof}
Equation (10.2) is (1.6) with (10.1).  Moving the normal argument from
\((k-np)/\sqrt{npq}\) to \((k+1/2-np)/\sqrt{npq}\) costs at most the
Lipschitz constant \(2/5\) times \(1/(2\sqrt{npq})\).
\end{proof}

The local theorem identifies individual binomial masses with the Gaussian
histogram.  Write \(\theta=k/n\),
\[
 D(\theta\Vert p)=
 \theta\log\frac{\theta}{p}+(1-\theta)\log\frac{1-\theta}{q},             \tag{10.3}
\]
and let \(m=\min(p,q)\).

\begin{theorem}[Quantitative local de Moivre--Laplace]
\label{thm:local-demoivre}
Fix a rational \(A>0\).  Suppose
\(|k-np|\le A\sqrt n\) and \(A/\sqrt n\le m/2\).  Define
\[
 E_n(A,p)=
 \frac{2A}{m\sqrt n}+
 \frac{4A^3}{3m^2\sqrt n}+
 \frac{1/12+1/(3m)}{n}.                               \tag{10.4}
\]
If \(E_n(A,p)\le1\), then
\[
 \Pp(B_n=k)
 =\frac{e^{-(k-np)^2/(2npq)}}{\sqrt{2\pi npq}}
   (1+\varepsilon_{n,k}),
 \qquad |\varepsilon_{n,k}|\le2E_n(A,p).             \tag{10.5}
\]
Consequently the piecewise-linear histogram in every fixed central window
converges uniformly to the normal density.
\end{theorem}

\begin{proof}
The constructive Robbins--Stirling bounds give
\[
 \log(r!)=(r+\tfrac12)\log r-r+\tfrac12\log(2\pi)+R_r,
 \qquad \frac1{12r+1}<R_r<\frac1{12r}.               \tag{10.6}
\]
Substitution in
\(\binom nkp^kq^{n-k}\) yields the exact representation
\[
 \Pp(B_n=k)=
 \frac{\exp[-nD(\theta\Vert p)+R_n-R_k-R_{n-k}]}
      {\sqrt{2\pi n\theta(1-\theta)}}.               \tag{10.7}
\]
For \(d=\theta-p\), Taylor's theorem at \(p\) gives
\[
 \left|D(p+d\Vert p)-\frac{d^2}{2pq}\right|
 \le\frac{4|d|^3}{3m^2},                             \tag{10.8}
\]
because throughout the interval between \(p\) and \(p+d\), both coordinates
are at least \(m/2\), and the absolute third derivative is at most
\(8/m^2\).  The logarithm of the prefactor ratio satisfies
\[
 \frac12\left|
 \log\frac{pq}{\theta(1-\theta)}
 \right|\le\frac{2|d|}{m}.                            \tag{10.9}
\]
Finally, \(k,n-k\ge nm/2\), so
\[
 |R_n-R_k-R_{n-k}|
 \le\frac{1/12+1/(3m)}{n}.                           \tag{10.10}
\]
Equations (10.8)--(10.10), together with \(|d|\le A/\sqrt n\), show that
the logarithm of the ratio between (10.7) and the Gaussian mass in (10.5)
has absolute value at most (10.4).  For \(|u|\le1\),
\(|e^u-1|\le2|u|\), proving (10.5).

All logarithms in this proof are finite regular-name computations with
Taylor remainder bounds.  The internal statement clears the positive
interval denominators and compares rational enclosures.
\end{proof}

\section{Local limits, densities, and Fourier formulations}
\label{sec:local-limit}

The ordinary CLT concerns CDFs.  Convergence of point masses or densities is
stronger and needs an aperiodicity or smoothing hypothesis.  For finite
integer laws one can give a fully constructive lattice theorem.

\begin{definition}[Lattice span]
Let \(X\) be an integer-valued finite law.  Its span is one if the greatest
common divisor of all differences between support points of positive weight
is one.  This is a bounded Euclidean-algorithm computation.
\end{definition}

\begin{theorem}[Effective lattice local CLT]\label{thm:lattice-local}
Let \(X\) be an integer-valued finite rational law of span one, mean \(\mu\),
variance \(\sigma^2>0\), and let \(S_n=X_1+\cdots+X_n\).  There is a
primitive-recursive modulus \(N(k)\) such that, for \(n\ge N(k)\),
\[
 \sup_{r\in\mathbb Z}
 \left|
  \sigma\sqrt n\,\Pp(S_n=r)
  -\phiN\!\left(\frac{r-n\mu}{\sigma\sqrt n}\right)
 \right|\le2^{-k}.                                    \tag{11.1}
\]
For span \(h\), the left mass is multiplied by \(\sigma\sqrt n/h\), with
only the admissible lattice points considered.
\end{theorem}

\begin{proof}
Let \(\psi(t)=\E e^{it(X-\mu)}\).  Integer Fourier inversion is the finite
orthogonality identity
\[
 \Pp(S_n=r)=\frac1{2\pi}\int_{-\pi}^{\pi}
    \psi(t)^n e^{-it(r-n\mu)}\dd t.                  \tag{11.2}
\]
It can be proved first for the finite Laurent polynomial \(\psi^n\), where
integration merely selects its zero coefficient.  After \(u=\sigma\sqrt n t\),
(11.2) becomes
\[
 \sigma\sqrt n\,\Pp(S_n=r)
 =\frac1{2\pi}\int_{-\pi\sigma\sqrt n}^{\pi\sigma\sqrt n}
 \psi\!\left(\frac{u}{\sigma\sqrt n}\right)^n
 e^{-iu(r-n\mu)/(\sigma\sqrt n)}\dd u.               \tag{11.3}
\]

Fix a rational \(0<\delta<\pi\).  On \(|t|\le\delta\), the independent-copy
argument gives \(|\psi(t)|\le e^{-ct^2}\) for an explicitly computed rational
\(c>0\).  Taylor's finite remainder gives, uniformly on every fixed
\(|u|\le A\),
\[
 \psi\!\left(\frac{u}{\sigma\sqrt n}\right)^n
 \longrightarrow e^{-u^2/2}                          \tag{11.4}
\]
with a primitive-recursive modulus.  On the intermediate region
\(A\le|u|\le\delta\sigma\sqrt n\), the damping bound makes both the
integrand and the Gaussian tail smaller than an explicit multiple of
\(e^{-c'A^2}\).

It remains to control \(\delta\le|t|\le\pi\).  Span one implies
\(|\psi(t)|<1\) there: equality in the triangle inequality would force all
phases \(e^{itx}\) on the support to coincide, hence \(t\) to be a multiple
of \(2\pi\).  Compactness is replaced by a finite rational mesh.  A known
Lipschitz constant for \(\psi\), together with exact interval evaluation at
mesh points, produces a rational \(r_\delta<1\) such that
\[
 \max_{\delta\le|t|\le\pi}|\psi(t)|\le r_\delta.      \tag{11.5}
\]
The outer integral is therefore at most
\(2\pi\sigma\sqrt n\,r_\delta^n\), with a bounded-search cutoff.

Finally, the inverse Fourier integral of \(e^{-u^2/2}\) is
\(\phiN(x)\).  The three regions just described give a uniform bound because
the phase in (11.3) has modulus one.  Choosing \(A\), then the mesh and
\(n\), yields (11.1).  Every choice is primitive recursive.
\end{proof}

For the Bernoulli law, \cref{thm:lattice-local} recovers the qualitative local
de Moivre--Laplace theorem; \cref{thm:local-demoivre} gives its sharper
central-window error.  For absolutely continuous laws, an analogous density
theorem follows from Fourier inversion when the input supplies an integrable
characteristic-function majorant.  Such a majorant is genuinely stronger
than a finite second moment, exactly as the source pages note.

\section{The multivariate central limit theorem}
\label{sec:multivariate}

A finite rational vector law in dimension \(d\) is a finite law with support
in \(\Q^d\).  Its mean \(\mu\) and covariance matrix
\(\Sigma=\E[(X-\mu)(X-\mu)^{\mathsf T}]\) are exact rational data.  For a
rational vector \(a\), write
\[
 \sigma_a^2=a^{\mathsf T}\Sigma a,
 \qquad \rho_a=\E|a^{\mathsf T}(X-\mu)|^3.           \tag{12.1}
\]

\begin{theorem}[Constructive Cram\'er--Wold CLT]\label{thm:cramer-wold}
Let \(X_1,X_2,\ldots\) be iid copies of a finite rational vector law.  For
every rational \(a\) with \(\sigma_a>0\),
\[
 \sup_x\left|
 \Pp\!\left[
  a^{\mathsf T}\frac{\sum_{i<n}(X_i-\mu)}{\sqrt n}
  \le x\right]
 -\PhiN(x/\sigma_a)
 \right|
 \le\frac{64\rho_a}{\sigma_a^3\sqrt n}.             \tag{12.2}
\]
If \(\sigma_a=0\), the projection is identically zero.  Consequently the
vector sums converge, in the rational-projection sense, to the multivariate
normal \(N_d(0,\Sigma)\).
\end{theorem}

\begin{proof}
The scalar variables \(a^{\mathsf T}(X_i-\mu)\) form an iid finite rational
law with variance and third moment (12.1).  Apply
\cref{cor:iid-rate}.  A multivariate normal is constructively specified by
requiring each rational projection to have the corresponding one-dimensional
normal law; consistency follows from the characteristic function
\[
 \phi_\Sigma(t)=\exp(-\tfrac12t^{\mathsf T}\Sigma t). \tag{12.3}
\]
Finite-dimensional characteristic functions are regular names, and the
quadratic form is rational at rational \(t\).
\end{proof}

For a fixed dimension, rational-projection convergence implies convergence
on every rational box whose boundary has zero \(N_d(0,\Sigma)\)-probability.
Here is the constructive content of the usual Cram\'er--Wold device.
Approximate a box indicator above and below by tensor products of
one-dimensional piecewise-linear ramps of width \(h\).  Each ramp is a finite
linear combination of positive-part functions and can be uniformly
approximated on a large rational cube by a trigonometric polynomial.  Its
expectation is therefore a finite linear combination of
characteristic-function values.  Gaussian tails control the complement of
the cube.  If coordinate \(i\) has variance \(\Sigma_{ii}>0\), its two
boundary layers have probability at most
\(2h/\sqrt{2\pi\Sigma_{ii}}\); if \(\Sigma_{ii}=0\), a continuity box is
required not to place a face on that deterministic coordinate.  Summing the
resulting regular-name bounds over the finitely many coordinates gives an
explicit modulus.  Fixing \(h\), the cube, and the trigonometric mesh before
the induction avoids compactness and choice.

The dimension-dependent convex-set estimate
\(C d^{1/4}\sum_i\E\|\Sigma^{-1/2}X_i\|^3\) stated on the source page is a
substantially deeper geometric Berry--Esseen theorem.  This paper proves all
rational projections with an explicit constant and fixed-dimensional box
convergence, but does not claim the optimal convex-set inequality.

\section{A random number of summands}
\label{sec:random-sum}

We now give an effective version of the Robbins random-sum theorem.  All laws
below are finite rational laws at each index, and every asserted convergence
comes with a primitive-recursive uniform-CDF modulus.

\begin{theorem}[Effective Robbins CLT]\label{thm:robbins}
Let \(X_1,X_2,\ldots\) be iid with mean \(\mu\), variance
\(\sigma^2\), and an effective finite-variance presentation.  Let \(N_n\) be
nonnegative integer-valued laws, independent of the \(X_i\), with
\[
 m_n=\E N_n,
 \qquad v_n=\Var(N_n)>0,
 \qquad \frac{N_n-m_n}{\sqrt{v_n}}\Longrightarrow N(0,1).               \tag{13.1}
\]
Assume \(\sigma^2+\mu^2>0\).  Then
\[
 \frac{\sum_{i=1}^{N_n}X_i-\mu m_n}
      {A_n}
 \Longrightarrow N(0,1),
 \qquad A_n^2=\sigma^2m_n+\mu^2v_n.                  \tag{13.2}
\]
The modulus is primitive recursive from the moduli in the hypotheses.
\end{theorem}

\begin{proof}
First, (13.1) forces \(v_n\to\infty\) and
\(m_n/\sqrt{v_n}\to\infty\).  For the second assertion, if the ratio stayed
below a rational \(C\), then the standardized nonnegative variable would
vanish below \(-C\), contradicting convergence to a normal CDF at a rational
point less than \(-C\).  For the first, Chebyshev puts at least \(3/4\) of
the mass of \(N_n\) in an interval containing at most
\(4\sqrt{v_n}+3\) integers.  If \(v_n\) stayed bounded, one atom would have a
fixed positive mass, impossible under uniform convergence to a continuous
CDF.  These contradictions have effective finite cutoffs.

Let \(\psi(u)=\E e^{iu(X-\mu)}\).  Conditional on \(N_n\), the
characteristic function of the numerator in (13.2) is
\[
 \E\left[
 e^{it\mu(N_n-m_n)/A_n}\psi(t/A_n)^{N_n}
 \right].                                             \tag{13.3}
\]
For a finite third-moment approximation to \(X\), (4.7) and a finite product
estimate give, for fixed \(t\),
\[
 \psi(t/A_n)^{N_n}
 =\exp\!\left(-\frac{\sigma^2t^2N_n}{2A_n^2}\right)+o(1)               \tag{13.4}
\]
in expectation.  Indeed, the expected error is bounded by a constant times
\(m_n/A_n^3\), which tends to zero when \(\sigma>0\).  The general effective
finite-variance case follows by the transfer argument of
\cref{thm:effective-iid}.  If \(\sigma=0\), the sum is \(\mu N_n\) and the
conclusion follows directly from (13.1).

Replace \(N_n\) by \(m_n\) in the real exponential in (13.4).  Its expected
change is bounded by
\[
 \frac{\sigma^2t^2}{2A_n^2}\E|N_n-m_n|
 \le\frac{\sigma^2t^2\sqrt{v_n}}{2A_n^2}\longrightarrow0.             \tag{13.5}
\]
When \(\mu=0\), this is \(t^2\sqrt{v_n}/(2m_n)\); otherwise the denominator
contains \(\mu^2v_n\).

Put
\[
 q_n=\frac{\sigma\sqrt{m_n}}{A_n},
 \qquad r_n=\frac{\mu\sqrt{v_n}}{A_n},
 \qquad q_n^2+r_n^2=1.                                \tag{13.6}
\]
After (13.4)--(13.5), expression (13.3) differs by \(o(1)\) from
\[
 e^{-q_n^2t^2/2}
 \E\exp\!\left(it r_n\frac{N_n-m_n}{\sqrt{v_n}}\right).               \tag{13.7}
\]
Uniform CDF convergence in (13.1), Chebyshev tails, and finite integration by
parts give convergence of the characteristic functions uniformly for
\(|r_n t|\le|t|\).  Thus (13.7) tends to
\(e^{-q_n^2t^2/2}e^{-r_n^2t^2/2}=e^{-t^2/2}\).
To use \cref{thm:esseen}, split its integral at a small dyadic \(\delta\).
The variable in (13.2) has exact variance one, so on \(|t|\le\delta\),
\(|\phi_n(t)-e^{-t^2/2}|/|t|\le |t|\); on
\(\delta\le|t|\le T\), the just-proved uniform characteristic-function
modulus applies.  Choose \(T\), then \(\delta\), then the row cutoff.  This
converts the compact-frequency estimate into a CDF modulus without hiding the
singularity at zero.  All error terms have explicit rational majorants, so
the proof arithmetizes.
\end{proof}

\section{Weak dependence and variance calculation}
\label{sec:dependence}

The full strong-mixing theorem on the source page is possible in an effective
presentation, but its mixing-coupling bookkeeping is much longer than the
independent proof.  We prove a substantial exact subclass: bounded stationary
finite-range dependence.  Such a process is strongly mixing with mixing
coefficient zero beyond its range.

\begin{definition}[Finite-range dependence]
A stationary finite rational process \((X_i)\) is \(m\)-dependent if finite
coordinate blocks whose index sets are separated by more than \(m\) are
independent.  Internally this is a family of exact finite-factorization
certificates.
\end{definition}

Let \(\E X_i=0\), \(|X_i|\le B\), and
\(\gamma_h=\E(X_0X_h)\).  Put
\[
 \sigma^2=\gamma_0+2\sum_{h=1}^{m}\gamma_h.           \tag{14.1}
\]

\begin{theorem}[Finite-range dependent CLT]\label{thm:m-dependent}
If \(\sigma^2>0\), then
\[
 \frac{X_1+\cdots+X_n}{\sigma\sqrt n}
 \Longrightarrow N(0,1)                              \tag{14.2}
\]
with a primitive-recursive modulus.  Moreover, for \(n>m\),
\[
 \Var(S_n)=n\sigma^2-2\sum_{h=1}^{m}h\gamma_h.        \tag{14.3}
\]
\end{theorem}

\begin{proof}
Expanding the square and grouping pairs by their lag gives
\[
 \Var(S_n)=n\gamma_0+2\sum_{h=1}^{m}(n-h)\gamma_h,   \tag{14.4}
\]
which is (14.3).

Choose a block length \(\ell=\lfloor n^{1/3}\rfloor\).  Partition the first
\(n\) coordinates into big blocks of length \(\ell\), separated by gaps of
length \(m\), plus a final remainder.  The big-block sums are independent.
The number of discarded coordinates is
\[
 O_m(n/\ell+\ell).                                    \tag{14.5}
\]
For any set \(A\) of coordinates, finite-range dependence and
\(|\Cov(X_i,X_j)|\le B^2\) give
\[
 \Var\!\left(\sum_{i\in A}X_i\right)
 \le(2m+1)B^2|A|.                                    \tag{14.6}
\]
Thus the discarded sum divided by \(\sqrt n\) tends to zero in probability
with an explicit Chebyshev modulus.

Each big-block sum is bounded by \(\ell B\).  Its total variance differs
from \(n\sigma^2\) by at most a rational constant times
\(n/\ell+\ell\), by (14.3), (14.5), and (14.6).  After normalizing its exact
total variance, \cref{thm:bounded-main} gives an error
\[
 O_{m,B,\sigma}(\ell/\sqrt n)=O(n^{-1/6}).            \tag{14.7}
\]
The relative variance error is \(O(n^{-1/3})\).  For the discarded sum,
choose a normalized threshold shift \(h=n^{-1/9}\): Chebyshev transfer
then contributes \(O(n^{-1/9})\), while normal-density continuity contributes
\(O(n^{-1/9})\) as well.  Every hidden constant is an explicit rational
expression in \(m,B,\sigma^{-1}\).  Bounded search turns these displayed
powers into a primitive-recursive cutoff, proving (14.2).
\end{proof}

\subsection{Weighted variance}

For \(g_i=g(i/n)\), stationarity gives the exact identity
\[
 \Var\!\left(\sum_{i=1}^ng_iX_i\right)
 =\gamma_0\sum_i g_i^2
  +2\sum_{h=1}^{n-1}\gamma_h\sum_{i=1}^{n-h}g_ig_{i+h}.                \tag{14.8}
\]
Assume a primitive-recursive modulus of continuity for \(g\), boundedness,
an effective absolutely summable covariance tail, and \(\gamma_n>0\) on the
indices under consideration.  For each fixed \(h\),
\(g_{i+h}-g_i\) is uniformly small, while the omitted boundary has only
\(h\) terms.  Split the lag sum at an effectively chosen \(H\), use the
covariance-tail modulus beyond \(H\), and obtain
\[
 \frac{\Var(\sum_i g_iX_i)}{n\gamma_n}\longrightarrow\sigma^2,
 \qquad \gamma_n=\frac1n\sum_i g_i^2.                \tag{14.9}
\]
This is the constructive content of the variance lemma stated on the source
page.  In the \(m\)-dependent case the tail is literally zero, and the same
blocking proof gives the weighted CLT whenever \(\gamma_n\) has a supplied
positive lower bound.

\section{Martingale differences on finite probability trees}
\label{sec:martingale}

A finite probability tree has rational transition weights.  A filtration is
the sequence of prefix partitions.  Conditional expectation is therefore a
finite weighted average at each node.  A martingale-difference row is a list
\(D_1,\ldots,D_r\) satisfying
\[
 \E[D_k\mid\mathcal F_{k-1}]=0                       \tag{15.1}
\]
at every node.  Put
\[
 v_k=\E[D_k^2\mid\mathcal F_{k-1}],
 \qquad V_j=\sum_{k\le j}v_k.                         \tag{15.2}
\]

\begin{theorem}[Bounded exact-variance martingale CLT]
\label{thm:bounded-martingale}
Suppose \(|D_k|\le b\) at every leaf and \(V_r=1\) at every leaf.  Then for
each rational \(T>0\),
\[
 \sup_{|t|\le T}
 \left|\E e^{it\sum_kD_k}-e^{-t^2/2}\right|
 \le e^{T^2b^2/2}
      \left(\frac{bT^3}{6}+\frac{b^2T^4}{8}\right). \tag{15.3}
\]
Consequently a primitive-recursive sequence with \(b_n\to0\) satisfies the
CLT, with an explicit modulus.  For example, taking
\(T=b^{-1/4}\) in Esseen smoothing gives a computable error tending to zero.
\end{theorem}

\begin{proof}
Let \(S_j=\sum_{k\le j}D_k\) and define the compensated process
\[
 H_j=\exp\!\left(itS_j+\frac{t^2}{2}(V_j-1)\right).   \tag{15.4}
\]
Because \(V_j\le V_r=1\), one has \(|H_j|\le1\).  Conditional Taylor
expansion and (15.1) give
\[
 \left|\E(e^{itD_k}\mid\mathcal F_{k-1})
        -1+\frac{t^2v_k}{2}\right|
 \le\frac{|t|^3}{6}\E(|D_k|^3\mid\mathcal F_{k-1})
 \le\frac{b|t|^3v_k}{6}.                             \tag{15.5}
\]
Also
\[
 \left|e^{-t^2v_k/2}-1+\frac{t^2v_k}{2}\right|
 \le\frac{t^4v_k^2}{8}.                              \tag{15.6}
\]
Thus, conditioning on \(\mathcal F_{k-1}\),
\[
 |\E(H_k\mid\mathcal F_{k-1})-H_{k-1}|
 \le e^{t^2v_k/2}
  \left(\frac{b|t|^3v_k}{6}+\frac{t^4v_k^2}{8}\right).                \tag{15.7}
\]
Now \(v_k\le b^2\), \(\sum v_k=1\), and
\(\sum v_k^2\le b^2\).  Telescope expectations from
\(H_0=e^{-t^2/2}\) to \(H_r=e^{itS_r}\).  This proves (15.3).
The smoothing integral on \([-T,T]\) is finite and its displayed bound is a
primitive-recursive expression.
\end{proof}

\begin{theorem}[Effective martingale Lindeberg theorem]
\label{thm:martingale-general}
Consider a primitive-recursive sequence of finite martingale-difference rows.
Assume effective moduli for
\[
 \Pp(|V_r-1|>\eta)\longrightarrow0                   \tag{15.8}
\]
for every dyadic \(\eta>0\), and
\[
 \sum_k\E\!
  \left[D_k^2\ind_{\{|D_k|>\eta\}}\right]
 \longrightarrow0.                                   \tag{15.9}
\]
Then \(\sum_kD_k\Longrightarrow N(0,1)\) with a primitive-recursive
modulus.
\end{theorem}

\begin{proof}
Fix the final precision and a small dyadic \(\eta\).  Truncate each difference
at \(\eta\) and subtract its conditional mean.  The sum of the discarded
parts has \(L^2\) norm bounded by an explicit multiple of (15.9), exactly as
in (9.3)--(9.4).  The truncated differences are bounded by \(2\eta\), and
their predictable variance differs from \(V_r\) by a quantity controlled by
(15.9).

On the good event from (15.8), stop just before the cumulative predictable
variance first exceeds one.  If a crossing would occur, the remaining
deficit is at most the next conditional variance, hence at most
\(4\eta^2\); if no crossing occurs, it is bounded by the supplied variance
error.  Split that rational deficit into finitely many rational pieces
\(r\le4\eta^2\).  For each piece append, conditionally on the current leaf, a
three-point increment taking \(2\eta,-2\eta,0\) with probabilities
\(r/(8\eta^2),r/(8\eta^2),1-r/(4\eta^2)\).  It has conditional mean zero,
variance exactly \(r\), and amplitude at most \(2\eta\).  Thus no irrational
branch value is introduced.  On each good branch the enlarged row now has
predictable variance exactly one and maximum increment tending to zero, so
\cref{thm:bounded-martingale} applies.  The stopped or modified bad branches
have total probability bounded by the supplied errors in (15.8)--(15.9).
Chebyshev transfer and normal anti-concentration then return to the original
row.

All stopping, truncation, branch extension, and conditional averaging occur
on a finite tree.  After the requested precision is fixed, the proof is a
bounded induction over nodes.  This is the constructive finite-tree version
of the martingale theorem stated on the source page.
\end{proof}

\section{Products, lognormal limits, and the law of large numbers}
\label{sec:products-lln}

\subsection{Products of positive variables}

Let \(X\) be a positive finite rational law.  Although \(\log X\) need not be
rational, its values have primitive-recursive regular names.  Put
\[
 \nu=\E\log X,
 \qquad \tau^2=\Var(\log X)>0,
 \qquad \rho_{\log}=\E|\log X-\nu|^3.                \tag{16.1}
\]
For \(P_n=\prod_{i=1}^nX_i\),
\[
 \frac{\log P_n-n\nu}{\tau\sqrt n}
 =\frac1{\tau\sqrt n}\sum_{i=1}^n(\log X_i-\nu).    \tag{16.2}
\]

\begin{theorem}[Multiplicative CLT]
For every positive rational threshold \(r\),
\[
 \left|
 \Pp(P_n\le r)-
 \PhiN\!\left(\frac{\log r-n\nu}{\tau\sqrt n}\right)
 \right|
 \le\frac{64\rho_{\log}}{\tau^3\sqrt n}.             \tag{16.3}
\]
The same conclusion under merely finite variance of \(\log X\) follows from
\cref{thm:effective-iid} when a logarithmic second-moment-tail modulus is
supplied.
\end{theorem}

\begin{proof}
The logarithm is strictly increasing, so \(P_n\le r\) is exactly
\(\sum_i\log X_i\le\log r\).  The product comparison is decidable because
both \(P_n\) and \(r\) are rational; only the Gaussian argument uses regular
names.  Apply \cref{cor:iid-rate}.  In an ordinary real completion, putting
\(r=e^{n\nu+\tau\sqrt n\,y}\) recovers the familiar lognormal-quantile
display at every continuity threshold.  The internal checker retains the
stronger proof-theoretically safe rational-threshold form above.
\end{proof}

Thus products become asymptotically lognormal after logarithmic centering and
scaling.  The essential moment condition is finite variance of \(\log X\),
not square-integrability of a density of \(X\).  A density assumption may be
a sufficient hypothesis in a particular local theorem, but it is not the
natural general condition for the multiplicative CLT.

\subsection{Weak law and the fluctuation scale}

For iid finite laws, the weak law of large numbers is already an exact finite
inequality:
\[
 \Pp(|\overline X_n-\mu|\ge\varepsilon)
 \le\frac{\sigma^2}{n\varepsilon^2}.                 \tag{16.4}
\]
It follows from Markov's inequality applied to
\((\overline X_n-\mu)^2\).  The LLN controls the order-one location of the
sample mean.  The CLT refines it by showing that the residual fluctuation has
order \(n^{-1/2}\):
\[
 S_n= n\mu+\sigma\sqrt n\,Z_n,
 \qquad Z_n\Longrightarrow N(0,1).                   \tag{16.5}
\]
This makes precise the asymptotic expansion described on the source page.

An almost-sure infinite-sequence statement is not literally a first-order
claim about one completed random sequence in \(\ipra\).  What is internal is
every finite maximal inequality.  For independent centered finite laws,
Kolmogorov's argument gives
\[
 \Pp\!\left(\max_{j\le n}|S_j|\ge a\right)
 \le\frac{\Var(S_n)}{a^2}.                            \tag{16.6}
\]
The proof partitions by the first crossing time and uses orthogonality of the
remaining centered sum.  Hence all finite certificates normally used in a
dyadic proof of the strong law are available; an actual infinite almost-sure
object belongs to a completion theory rather than to the first-order core.

\section{Stable laws and the generalized theorem}
\label{sec:stable}

The normal law is stable: sums of independent
\(N(\mu,\sigma^2)\) variables are exactly normal, as their characteristic
functions multiply.  The Cauchy law is also stable, because its characteristic
function is \(e^{-|t|}\) and
\((e^{-|t|})^n=e^{-n|t|}\).  These identities have constructive regular-name
proofs by finite integration with explicit tails.

The difficult direction of the generalized central limit theorem says that
every nondegenerate limit of centered and scaled iid sums is stable.  A useful
constructive form separates the compactness step from the algebraic step.

\begin{proposition}[Effective stability criterion]
Suppose a fixed iid law and primitive-recursive sequences \(a_n>0,b_n\)
satisfy
\[
 a_n(X_1+\cdots+X_n)-b_n\Longrightarrow Z            \tag{17.1}
\]
with a modulus.  Assume additionally that, for every fixed \(m\), regular
names and moduli are supplied for
\[
 c_m=\lim_n\frac{a_{mn}}{a_n},
 \qquad
 d_m=\lim_n\left(\frac{a_{mn}}{a_n}m b_n-b_{mn}\right).                \tag{17.2}
\]
Then, for independent copies \(Z_1,\ldots,Z_m\),
\[
 Z\ \stackrel d=\ c_m(Z_1+\cdots+Z_m)+d_m.            \tag{17.3}
\]
Thus \(Z\) is stable.
\end{proposition}

\begin{proof}
Partition the sum of \(mn\) variables into \(m\) independent blocks of
length \(n\).  Rewrite its normalization using (17.2).  The joint
characteristic function of the normalized blocks is the product of their
one-block characteristic functions.  Pass to the supplied limits with the
explicit moduli to obtain (17.3).
\end{proof}

The unrestricted classical theorem must prove that suitable ratio limits in
(17.2) exist, usually by subsequence compactness and regular-variation
arguments.  Those qualitative existence steps do not automatically provide
primitive-recursive moduli.  They are therefore recorded as an open
completion layer, not asserted as an internal \(\ipra\) theorem in this
package.

\section{Counterexamples and misconceptions}
\label{sec:counterexamples}

\subsection{No universal ``sample size 30'' rule}

Let
\[
 \Pp(X=100)=\Pp(X=-100)=\frac1{20000},
 \qquad \Pp(X=0)=\frac{9999}{10000}.                  \tag{18.1}
\]
Then \(\E X=0\) and \(\Var X=1\).  For \(n=30\),
\[
 \Pp(S_{30}=0)\ge\left(\frac{9999}{10000}\right)^{30}
 \ge1-\frac{30}{10000}=0.997.                        \tag{18.2}
\]
Any distribution with an atom of mass \(p\) has Kolmogorov distance at least
\(p/2\) from every continuous distribution: the left and right CDF limits
differ by \(p\), while the comparison CDF does not jump.  Hence the
standardized sum in (18.2) is more than \(0.498\) away from normal.  This is
an exact rational disproof of a distribution-free ``\(n\ge30\)'' heuristic.

\subsection{Pairwise independence is not enough}

Let \(U\) be uniform on \(\mathbb F_2^m\).  For every nonzero
\(a\in\mathbb F_2^m\), put
\[
 X_a=(-1)^{a\cdot U}.                                 \tag{18.3}
\]
The \(n=2^m-1\) variables are Rademacher and pairwise independent, because
two distinct nonzero linear functionals over \(\mathbb F_2\) are jointly
uniform.  Yet
\[
 \sum_{a\ne0}X_a=
 \begin{cases}
 n,&U=0,\\
 -1,&U\ne0.
 \end{cases}                                          \tag{18.4}
\]
Thus the variance-normalized sum converges in probability to zero, not to a
normal law.  Full independence, or a genuine weak-dependence hypothesis, is
essential.

\subsection{Sampling does not Gaussianize the population}

The empirical distribution of individual observations approximates the
population law, not a Gaussian law.  The CLT concerns the separate random
variable formed by a sum or average across repeated samples.  This logical
distinction is visible in the finite coding: the empirical histogram is a
vector of category counts, whereas the sample mean is one rational observable
of the product space.

\section{Examples and elementary applications}
\label{sec:applications}

\begin{example}[Fair dice]
For a fair die \(X\in\{1,\ldots,6\}\),
\[
 \mu=\frac72,
 \qquad \sigma^2=\frac{35}{12},
 \qquad \rho_3=\frac{51}{8}.                         \tag{19.1}
\]
Therefore the normalized sum of \(n\) dice has Kolmogorov error at most
\[
 \min\!\left(1,
 \frac{64(51/8)}{(35/12)^{3/2}\sqrt n}\right).       \tag{19.2}
\]
Although this constant is intentionally coarse, (1.7) supplies a wholly
rational certified sample-size cutoff for any requested dyadic accuracy.
\end{example}

\begin{example}[Linear regression with fixed design]
In a fixed-design scalar regression, the centered least-squares estimator is
\[
 \widehat\beta-\beta
 =\frac{\sum_i w_i\varepsilon_i}{\sum_i w_i^2}.       \tag{19.3}
\]
For independent centered errors of variance \(\sigma^2\), the standardized
numerator is a triangular array.  If the errors are bounded by \(B\), then
\cref{thm:bounded-main} gives
\[
 \dk\le
 7B\frac{\max_i|w_i|}{\sigma\sqrt{\sum_iw_i^2}}.     \tag{19.4}
\]
Thus the usual ``no dominating observation'' condition
\(\max|w_i|/\sqrt{\sum w_i^2}\to0\) has an immediate constructive normal
approximation.  Third-moment errors can instead use \cref{thm:be-main}.
\end{example}

The same weighted-array argument applies to measurement errors, survey
averages, and other ``many small unrelated effects'' examples.  The phrase is
not a theorem until ``small'' and ``unrelated'' are replaced by a bound such
as (9.1), (14.6), or (19.4).

\section{Why the proofs fit intuitionistic primitive recursive arithmetic}
\label{sec:ipra-audit}

\subsection{The exact first-order theorem}

The mathematical notation in \cref{thm:be-main} abbreviates a universal
arithmetic sentence.  A natural number \(e\) codes a finite row, a natural
number \(u\) codes a signed rational threshold, and \(p\) is a requested
precision.  A fixed primitive-recursive checker
\[
 \mathsf{BECheck}(e,u,p)                               \tag{20.1}
\]
decodes the row, computes exact means, centered support values, total
variance and third moments, enumerates the product-space CDF at the raw
threshold, and evaluates interval names for the Gaussian CDF and the error
bound.  It returns zero exactly when the dyadically slackened enclosure
required by (1.4) succeeds.  Invalid encodings return zero.  The internal
open theorem is simply
\[
 \mathsf{BECheck}(e,u,p)=0,                            \tag{20.2}
\]
and the closed theorem universally quantifies \(p,u,e\).

Returning zero on invalid inputs is convenient: validity need not be placed
in an implication with logical connectives.  The checker can branch on its
decidable validity test and make (20.2) automatically true in the invalid
branch.  In the valid branch it checks integer inequalities obtained by
clearing explicitly positive denominators.

\subsection{Where induction occurs}

Every induction used in the proof has a quantifier-free invariant after all
external parameters are fixed:

\begin{center}
\begin{tabular}{p{0.31\textwidth}p{0.58\textwidth}}
\toprule
Construction & Quantifier-free invariant \\
\midrule
Finite sums and products & defining integer/rational equality after a fixed number of terms \\
Mixed-radix product space & decoded prefix has the stated accumulated weight and value \\
Expectation factorization & partial double sum equals the product of partial sums \\
Characteristic-product telescoping & partial product difference equals the corresponding partial replacement sum \\
Taylor evaluators & finite polynomial recurrence and a cleared-denominator remainder inequality \\
Bisection for roots & the current rational interval contains the target and has the claimed width \\
Riemann sums & partial sum plus an explicit rational error budget \\
Esseen integral & mesh prefix satisfies its accumulated upper/lower enclosure \\
Threshold CDF & accumulated weight equals the weight of all decoded outcomes satisfying a decidable comparison \\
\bottomrule
\end{tabular}
\end{center}

The proof never uses induction on a formula containing an unbounded
quantifier.  Universal closure is applied only after the open QF derivation
has been constructed.

\subsection{Intuitionistic logic}

No excluded-middle principle for arbitrary propositions occurs.  Case splits
are on decidable comparisons of integers or rationals, finite membership
predicates, or Boolean validity checks.  Positivity needed for division is
carried as a witness.  Existence claims such as ``choose a mesh'' use a
displayed primitive-recursive bound, not a negative double-negation
argument.  The independent copy, product space, and martingale tree are
finite constructions rather than uses of a countable product theorem.

The external paper occasionally writes a supremum, an improper integral, or
a real limit for readability.  The internal forms are respectively:

\begin{itemize}
\item a universal inequality at every rational threshold;
\item a compact rational Riemann sum plus a displayed tail bound;
\item a primitive-recursive modulus and a dyadic interval inclusion.
\end{itemize}

\subsection{The theorem-schema issue}

First-order \(\ipra\) has number variables, not variables ranging over all
primitive-recursive programs or all real functions.  Therefore the result is
a uniform metatheorem: for each fixed finite-law evaluator or fixed effective
presentation, there is a mechanically generated arithmetic derivation.
Equivalently, one may work in a definitional extension containing the fixed
primitive-recursive function symbols.  This is the same honest reading used
throughout the ProveIt constructive-analysis packets.

\section{Lean formalization blueprint}
\label{sec:lean}

The archive contains three Lean-facing artifacts:
\begin{itemize}[nosep]
\item \code{FiniteCLTBlueprint.lean};
\item \code{ObjectLanguageTarget.lean};
\item \code{ROADMAP.md}.
\end{itemize}
They are deliberately split into an ambient finite mathematics layer and a
deep object-language layer.

\subsection{Ambient finite mathematics}

A first module should define
\begin{verbatim}
structure FinLaw where
  card      : Nat
  value     : Fin card -> Rat
  weight    : Fin card -> Nat
  total     : Nat
  total_pos : 0 < total
  weight_sum : (sum i, weight i) = total
\end{verbatim}
with exact functions
\code{expect}, \code{mean}, \code{variance}, \code{thirdAbs},
\code{cdf}, \code{product}, and \code{convolution}.  Product outcomes use a
mixed-radix code.  The principal finite lemmas are:
\begin{verbatim}
expect_product
variance_sum_centered
charFun_product
product_replace_telescoping
cdf_sum_raw_threshold
\end{verbatim}
All are theorems about \code{Rat} or finite sums and can be proved before any
real analysis is introduced.

The analytic interface is a structure of regular rational interval names:
\begin{verbatim}
structure Regular where
  lo hi    : Nat -> Rat
  ordered : forall p, lo p <= hi p
  nestedLo : Monotone lo
  nestedHi : Antitone hi
  width    : forall p, hi p - lo p <= 1 / 2^p
\end{verbatim}
Separate modules then implement square root, sine, cosine, exponential,
\(\pi\), compact integration, Gaussian density/CDF, and complex pairs.  The
implementation should reuse the existing constructive calculus, Taylor, and
Stirling packets where their APIs are stable, but the CLT theorem must depend
only on finite evaluators and proved error bounds, not on semantic
\code{Real} completeness.

The analytic theorem order should mirror the paper:
\begin{enumerate}
\item finite Taylor inequalities (4.7)--(4.9);
\item Gaussian integral and characteristic function;
\item the kernel (5.1), its normalization, transform, and tail;
\item Esseen smoothing;
\item independent-copy damping and one-factor replacement;
\item bounded-array constant 7;
\item Berry--Esseen constant 64;
\item Lindeberg truncation and the corollaries.
\end{enumerate}

An ambient Lean theorem may use \code{Real} as a semantic convenience for
initial experimentation.  It is not the final endpoint.

\subsection{Object-language endpoint}

The decisive module imports
\code{ConciseIPRA.NativeBiInterpretation}.  It must build concrete
\code{PR 3} syntax for (20.1) and then construct
\begin{verbatim}
structure InternalProof extends Representation where
  derived : ([] : List (Fm PTm 3))
              |- [PRAAx] berryEsseenOpen check
\end{verbatim}
where the open formula is the equation
\code{check(row,q,p)=0}.  The archive's exact notation uses the repository's
sequent symbol rather than the ASCII rendering above.  The completed value
must be produced from explicit \code{PRAAx} derivations.  A semantic theorem
that the checker always returns zero is not a reflection principle and does
not fill \code{derived}.

After the open proof, universal generalization constructs the closed sentence.
The existing repository theorem \code{pra_iff_native} then transports that
proof to the native scalar calculus, and \code{FromPRA.code} compiles the
checker syntax to the concise iterator language.  Thus the final chain is
\[
 \text{finite algorithms}
 \longrightarrow \PR\text{ syntax}
 \longrightarrow \mathsf{PRAAx}\text{ derivation}
 \longrightarrow \text{native ConciseIPRA proof}.    \tag{21.1}
\]

\subsection{Concrete implementation order}

The lowest-risk route to completion is:

\begin{enumerate}
\item Compile signed integers, positive-denominator rationals, bounded sums,
products, powers, pairing, lists, and mixed-radix decoding to \code{PR}.
\item Prove the defining equations of the exact probability layer inside
\code{PRAAx}.
\item Implement regular names with fixed precision schedules.  Keep all
analytic algorithms closed and named; do not attempt a universal program
evaluator.
\item Formalize the finite Taylor and Gaussian lemmas.  Reuse previously
verified finite-integration components where possible.
\item Formalize Esseen smoothing first as a rational mesh theorem.  The
external integral notation should be a corollary of the mesh enclosure.
\item Construct the bounded theorem; it is simpler and tests all major APIs.
\item Add truncated damping and construct the constant-64 theorem.
\item Instantiate \code{Representation}, then \code{InternalProof}; close and
transport the theorem.
\item Put ordinary measure-theoretic interpretations and Mathlib probability
corollaries in downstream modules so they cannot mask a missing internal
proof.
\end{enumerate}

The repository revision inspected for this plan begins
\code{979613af}; the full SHA is recorded in \code{provenance/SOURCES.md}.
Its toolchain file specifies Lean \code{v4.32.0}.  The supplied Lean files contain no \code{sorry} and
introduce no central-limit axiom.  They were not compiled in the preparation
environment, because no Lean/Lake executable was installed.  Successful
elaboration is therefore not claimed.  The concrete unfinished endpoint is a
closed checker \code{beTest} together with an actual value of
\code{CentralLimitIPRAPlan.InternalProof beTest}.

\section{Coverage summary and remaining boundaries}
\label{sec:coverage}

The following compact table supplements the machine-readable-style ledger in
\code{THEOREM_LEDGER.md}.

\begin{center}
\small
\begin{longtable}{p{0.31\textwidth}p{0.16\textwidth}p{0.44\textwidth}}
\toprule
Source item & Status & Result here \\
\midrule
\endfirsthead
\toprule
Source item & Status & Result here \\
\midrule
\endhead
Classical iid CLT & effective/direct & Direct for finite rational laws; effective completion with a second-moment-tail/quantization modulus. \\
Pointwise and uniform CDF convergence & proved & Explicit universal-threshold Berry--Esseen inequality. \\
Berry--Esseen & proved & Independent triangular arrays, constant 64. \\
Lyapunov and Lindeberg & proved/effective & Quantitative truncation, all fixed rational \(\delta>0\), explicit nonconverse. \\
de Moivre--Laplace & proved & Global CDF rate and local Robbins--Stirling theorem. \\
Density/convolution form & proved/effective & Exact convolution and span-one lattice local CLT; continuous density form under a Fourier majorant. \\
Multivariate CLT & proved/effective & Every rational projection with rate; fixed-dimensional Gaussian-continuity boxes. \\
Convex-set \(d^{1/4}\) rate & open here & Requires separate high-dimensional geometry. \\
Random number of summands & effective & Constructive Robbins theorem with moduli. \\
Weak dependence & substantial subclass & Bounded stationary finite-range dependence; exact covariance formula. \\
Martingale differences & effective & Finite-tree predictable-variance/Lindeberg theorem. \\
Products/lognormality & proved/effective & Finite variance of \(\log X\), with a third-moment rate when available. \\
Generalized stable-law characterization & partial & Stability algebra and an effective ratio-limit criterion; compactness classification not reconstructed. \\
Entropy monotonicity & open here & Needs differential entropy and entropy-power machinery plus regularity. \\
Convex bodies, lacunary series, Gaussian polytopes, Haar matrices, subsequences, crystal lattices & open here & Distinct geometric, harmonic, or compactness theories; not consequences of the finite-array engine alone. \\
\bottomrule
\end{longtable}
\end{center}

The status word ``effective'' is not a euphemism for a missing proof.  It
means that the classical qualitative hypothesis has been replaced by the
explicit modulus needed for a first-order constructive theorem.  By contrast,
items marked open require additional mathematical machinery not developed in
this packet.

\appendix

\section{Rational audit of the two main constants}
\label{app:constants}

The bounded theorem uses only
\[
 \pi>\frac{25}{8},
 \qquad \frac1{\sqrt\pi}<\frac35,
 \qquad M=\frac1{\sqrt{2\pi}}<\frac25,
 \qquad \frac1\pi<\frac8{25}.                         \tag{A.1}
\]
Consequently
\[
 \frac{11}{12\sqrt\pi}<\frac{11}{20},
 \qquad
 \frac{48M}{\pi}<48\cdot\frac25\cdot\frac8{25}
 =\frac{768}{125},                                    \tag{A.2}
\]
and
\[
 \frac{11}{20}+\frac{768}{125}
 =\frac{3347}{500}<7.                                 \tag{A.3}
\]

For the Berry--Esseen theorem, \(5\pi<16\) and \(\pi>3\) imply
\[
 \frac{2}{\pi}
 \left(\frac{5\sqrt{5\pi}}{24}+\frac{25}{16}\right)
 <\frac23\left(\frac56+\frac{25}{16}\right)
 =\frac{115}{72}<\frac85.                             \tag{A.4}
\]
The smoothing term satisfies
\[
 \frac{384M}{\pi}
 <384\cdot\frac25\cdot\frac8{25}
 =\frac{6144}{125}.                                   \tag{A.5}
\]
Therefore
\[
 \frac85+\frac{6144}{125}
 =\frac{6344}{125}<64.                                \tag{A.6}
\]
Every comparison in (A.1)--(A.6) becomes an integer inequality after
squaring only nonnegative quantities.  The accompanying Python validation
records the exact \code{Fraction} calculations; floating-point values are
display-only.

\section{A checker-level statement}
\label{app:checker}

One convenient checker specification is the following.  Given
\((e,u,p)\):

\begin{enumerate}
\item Decode \(e\) as a row of finite rational law codes.  On failure, return
zero.
\item Compute exact centered values, \(s^2\), \(\rho\), and the raw-sum CDF
\(F(q)\), where \(q\) is decoded from \(u\).  If \(s^2=0\), check the
separate deterministic statement.
\item Request precision \(p+c(e,u)\) from the square-root, reciprocal,
Gaussian-CDF, and multiplication interval evaluators, where the guard
\(c\) is a displayed bit-length polynomial large enough for propagated
rounding errors.
\item Form rational intervals \([L_F,U_F]\), \([L_\Phi,U_\Phi]\), and
\([L_E,U_E]\) for the exact CDF, \(\PhiN(q/s)\), and
\(\min(1,64\rho/s^3)\).
\item Return zero when
\[
 U_F-L_\Phi\le U_E+2^{-p}
 \quad\text{and}\quad
 U_\Phi-L_F\le U_E+2^{-p};                            \tag{B.1}
\]
otherwise return one.
\end{enumerate}

Regularity makes the checker eventually accept at every precision, while the
mathematical proof establishes that the fixed precision schedule already
suffices.  Formula (B.1) contains only rational comparisons, hence integer
comparisons after clearing positive denominators.

\section{Further research questions}

\begin{enumerate}
\item Can the elementary finite Fourier proof be optimized inside \(\ipra\)
to a Berry--Esseen constant below one without materially enlarging the
checker?
\item Is there a proof-friendly multidimensional smoothing kernel giving a
reasonable explicit convex-set bound while retaining only QF induction?
\item Can the strong-mixing theorem with the source page's polynomial mixing
rate and twelfth moment be reconstructed with a compact primitive-recursive
block-coupling library?
\item What is the smallest regular-name and finite-integration core needed to
formalize Esseen smoothing in the native ConciseIPRA calculus?
\item Can the effective ratio hypotheses in \cref{sec:stable} be derived from
an intrinsic computable tightness condition, yielding a useful constructive
generalized CLT?
\item Can the finite-tree martingale proof be sharpened to a transparent
Berry--Esseen rate in terms of predictable third moments and variance error?
\item Which local limit hypotheses admit the shortest object-language
checker: lattice spectral gaps, integrable characteristic functions, or
bounded-variation densities?
\end{enumerate}

\begin{thebibliography}{9}

\bibitem{WikipediaCLT}
Wikipedia contributors,
\emph{Central limit theorem}, revision consulted 15 November 2026,
\url{https://en.wikipedia.org/wiki/Central_limit_theorem}.

\bibitem{MathWorldCLT}
E. W. Weisstein,
\emph{Central Limit Theorem}, MathWorld---A Wolfram Resource,
consulted 15 November 2026,
\url{https://mathworld.wolfram.com/CentralLimitTheorem.html}.

\bibitem{Feller}
W. Feller,
\emph{An Introduction to Probability Theory and Its Applications},
Volumes I--II, Wiley.

\bibitem{Petrov}
V. V. Petrov,
\emph{Sums of Independent Random Variables}, Springer.

\bibitem{Billingsley}
P. Billingsley,
\emph{Probability and Measure}, Wiley.

\bibitem{ProveIt}
V. Reshetnikov,
\emph{ProveIt}, repository revision \code{979613af}\allowbreak\ldots\code{d749f},
\url{https://github.com/VladimirReshetnikov/ProveIt}.

\end{thebibliography}

\end{document}
